\chapter{Conclusione}\label{chapter:conclusione}%Se si cambia il Titolo cambiare anche la riga successiva così che appia corretto nell'conclusione
%\addcontentsline{toc}{chapter}{Conclusione} %Per far apparire Introduzione nell'indice (Il nome deve rispecchiare quello del chapter)
Il lavoro svolto in questa tesi ha avuto come obiettivo 
la progettazione e l'implementazione di una metodologia computazionale 
per lo studio e il riconoscimento di potenziali 
\textit{clustering alternativi} nei dati di trascrittomica spaziale.  

\medskip
L'intero processo è stato sviluppato in ambiente Python, integrando librerie come \texttt{Scanpy}, \texttt{sklearn} e strumenti per l'analisi grafica e statistica, fino a includere modelli avanzati come \textit{GraphST} e \textit{Mclust}.  
In particolare, la pipeline ha previsto:  
\begin{itemize}
  \item la costruzione di una mappa cartesiana degli \textit{spot} e la definizione della lista degli spot vicini mediante l'algoritmo dei \textit{Nearest Neighbors};
  \item la suddivisione dei geni in sottogeni attraverso l'individuazione delle componenti connesse via \textit{Breadth-First Search};
  \item il calcolo dei rapporti spaziali tra sottogeni, classificati in cinque categorie (\texttt{DISGIUNTI}, \texttt{IN\_COSTA}, \texttt{A\_CONTENUTO}, \texttt{B\_CONTENUTO}, \texttt{VICINI});
  \item la rappresentazione dei rapporti tramite un grafo esportato e visualizzato in \textit{Gephi}, utile per l'analisi visiva dei pattern e delle connessioni tra geni;
\end{itemize}

\medskip
L'analisi del grafo dei geni \texttt{VICINI} ha rivelato la presenza di cluster ben distinti: un macro-cluster centrale, caratterizzato da un'elevata connettività, e uno secondario di dimensioni ridotte, in cui un singolo gene fungeva da nodo centrale.  
Questa osservazione ha confermato che i rapporti di vicinanza possono essere interpretati come indicatori di co-localizzazione spaziale tra geni, suggerendo possibili interazioni locali o domini funzionali condivisi.

\medskip
Per valutare se tali correlazioni spaziali corrispondessero anche a una reale struttura trascrittomica interna al tessuto, si è proceduto con un'analisi di clustering basata su \textit{GraphST} e \textit{Mclust}.  
Nonostante la solidità metodologica dell'approccio, i risultati non hanno mostrato l'emergere di cluster alternativi significativi.  
Questo esito suggerisce che la vicinanza spaziale tra geni non implica necessariamente una correlazione diretta di espressione, ma può riflettere fenomeni locali di prossimità o di contatto marginale.

\medskip
Il lavoro ha comunque raggiunto l'obiettivo di costruire una pipeline modulare, robusta e replicabile per l'analisi computazionale di dati di trascrittomica spaziale, aprendo la strada a possibili sviluppi futuri.  
Tra le direzioni di approfondimento più promettenti si segnalano:
\begin{itemize}
  \item l'integrazione con algoritmi di ricerca per scovare in autonomia potenziali nuovi pattern nel grafo Gephi;
  \item la creazione di un algoritmo di classificazione migliore basato sulla struttura delle maschere binarie;
\end{itemize}

\medskip
In conclusione, la metodologia proposta rappresenta un contributo utile per l'esplorazione delle relazioni spaziali tra geni nei dati di trascrittomica.  
Pur non avendo individuato nuovi cluster, il lavoro ha permesso di consolidare una base analitica solida e riproducibile, capace di essere estesa e adattata a futuri studi in ambito bioinformatico e computazionale.  
La trascrittomica spaziale resta un campo in rapida evoluzione, e strumenti di analisi come quello sviluppato in questa tesi potranno costituire un supporto prezioso per comprendere meglio la complessità dei tessuti e l'organizzazione spaziale dell'espressione genica.
